%% kapa: exam=2 round=36 year=2025 period=1 subject=appraisal_practice subject_ko=감정평가실무 seq=3 number=3 points=20
%% source: Q-Net 감정평가사 기출문제 (cst00309) artlSeq=5239325 file=raw/1008/5239325/2025년 제36회 감정평가사 2차 감정평가실무.pdf
%% conversion: kordoc 4.24.51 -> tools/gen2.py (auto); review=proofread-gpt-sol; flags=none
\begin{longquestion}{【문제 3】}
감정평가사 甲은 법원으로 부터 피고가 시행한 도시계획시설사업(이하 “공익사업”이라함)에 편입된 토지에 대하여 원고의 환매권 상실로 인한 손해액을 증명하기 위한 감정평가를 의뢰받았다. 관련 법규 및 이론을 참고하고, 제시된 자료를 활용하여 다음 물음에 답하시오. (20점) \par
\medskip\noindent 물음 1) 환매권 상실 당시를 기준으로 한 목적물의 감정평가액을 구하시오. (10점) \par
\medskip\noindent 물음 2) 손해액 산정을 위한 인근 유사 토지의 지가변동률을 구하시오. (10점) \par
< 자료 1 > 감정평가 대상의 확정 \par
{\small\centering
\resizebox{\ifdim\width>\linewidth\linewidth\else\width\fi}{!}{%
\begin{tabular}{|c|c|c|c|c|}
\hline
\makecell{일련번호} & \makecell{소재지} & \makecell{지번} & \makecell{면적(m²)} & \makecell{비고} \\
\hline
\makecell{1} & \makecell{A시 B구 C동} & \makecell{100} & \makecell{700} & \makecell{토지} \\
\hline
\makecell{2} & \makecell{A시 B구 C동} & \makecell{200} & \makecell{900} & \makecell{토지} \\
\hline
\end{tabular}}\par}

< 자료 2 > 법원 제시 기준일자 \par
{\small\centering
\resizebox{\ifdim\width>\linewidth\linewidth\else\width\fi}{!}{%
\begin{tabular}{|c|c|c|c|}
\hline
\makecell{일련번호} & \makecell{소유권 상실일} & \makecell{환매권 발생일} & \makecell{환매권 상실일} \\
\hline
\makecell{1} & \makecell{2008.12.29.} & \makecell{2013.12.29.} & \makecell{2014.12.29.} \\
\hline
\makecell{2} & \makecell{2016.06.17.} & \makecell{2021.06.17.} & \makecell{2022.06.17.} \\
\hline
\end{tabular}}\par}

< 자료 3 > 토지의 개황 \par
1. 소유권 상실일 당시 \par
{\small\centering
\resizebox{\ifdim\width>\linewidth\linewidth\else\width\fi}{!}{%
\begin{tabular}{|c|c|c|c|c|c|c|c|}
\hline
\makecell{일련번호} & \makecell{소재지} & \makecell{지번} & \makecell{지목} & \makecell{이용상황} & \makecell{용도지역} & \makecell{형상\\지세} & \makecell{도로접면} \\
\hline
\makecell{1} & \makecell{A시\\B구 C동} & \makecell{100} & \makecell{전} & \makecell{주거기타} & \makecell{자연녹지} & \makecell{사다리형\\평지} & \makecell{세로(가)} \\
\hline
\makecell{2} & \makecell{A시\\B구 C동} & \makecell{200} & \makecell{전} & \makecell{주거나지} & \makecell{2종일주} & \makecell{가장형\\평지} & \makecell{소로한면} \\
\hline
\end{tabular}}\par}

2. 환매권 발생일 및 환매권 상실일 당시 \par
{\small\centering
\resizebox{\ifdim\width>\linewidth\linewidth\else\width\fi}{!}{%
\begin{tabular}{|c|c|c|c|c|c|c|c|}
\hline
\makecell{일련번호} & \makecell{소재지} & \makecell{지번} & \makecell{지목} & \makecell{이용상황} & \makecell{용도지역} & \makecell{형상\\지세} & \makecell{도로접면} \\
\hline
\makecell{1} & \makecell{A시\\B구 C동} & \makecell{100} & \makecell{전} & \makecell{주거기타} & \makecell{자연녹지} & \makecell{부정형\\평지} & \makecell{소로한면} \\
\hline
\makecell{2} & \makecell{A시\\B구 C동} & \makecell{200} & \makecell{전} & \makecell{주거나지} & \makecell{2종일주} & \makecell{부정형\\평지} & \makecell{중로한면} \\
\hline
\end{tabular}}\par}

3. 당해 공익사업으로 인해 토지 등의 가격이 변동된 것으로 조사됨 \par
< 자료 4 > 표준지공시지가 \par
1. 인근 표준지공시지가 특성항목 \par
{\small\centering
\resizebox{\ifdim\width>\linewidth\linewidth\else\width\fi}{!}{%
\begin{tabular}{|c|c|c|c|c|c|c|c|}
\hline
\makecell{기호} & \makecell{소재지} & \makecell{지목} & \makecell{면적(m²)} & \makecell{이용상황} & \makecell{용도지역} & \makecell{형상\\지세} & \makecell{도로접면} \\
\hline
\makecell{1} & \makecell{A시 B구\\C동 60} & \makecell{전} & \makecell{700} & \makecell{전} & \makecell{자연녹지} & \makecell{부정형\\평지} & \makecell{맹지} \\
\hline
\makecell{2} & \makecell{A시 B구\\C동 80} & \makecell{전} & \makecell{500} & \makecell{전} & \makecell{자연녹지} & \makecell{부정형\\평지} & \makecell{세로(가)} \\
\hline
\makecell{3} & \makecell{A시 B구\\C동 120} & \makecell{전} & \makecell{400} & \makecell{주거기타} & \makecell{자연녹지} & \makecell{부정형\\평지} & \makecell{세로(가)} \\
\hline
\makecell{4} & \makecell{A시 B구\\C동 140} & \makecell{전} & \makecell{600} & \makecell{주거기타} & \makecell{자연녹지} & \makecell{부정형\\평지} & \makecell{소로한면} \\
\hline
\makecell{5} & \makecell{A시 B구\\C동 160} & \makecell{전} & \makecell{700} & \makecell{주거나지} & \makecell{2종일주} & \makecell{부정형\\완경사} & \makecell{세로(불)} \\
\hline
\makecell{6} & \makecell{A시 B구\\C동 180} & \makecell{전} & \makecell{600} & \makecell{주거나지} & \makecell{3종일주} & \makecell{가장형\\평지} & \makecell{세로(가)} \\
\hline
\makecell{7} & \makecell{A시 B구\\C동 220} & \makecell{전} & \makecell{900} & \makecell{주거기타} & \makecell{2종일주} & \makecell{부정형\\평지} & \makecell{세로(가)} \\
\hline
\makecell{8} & \makecell{A시 B구\\C동 240} & \makecell{전} & \makecell{800} & \makecell{주거나지} & \makecell{2종일주} & \makecell{부정형\\평지} & \makecell{소로한면} \\
\hline
\end{tabular}}\par}

주1) 기호 2, 4, 6, 8 표준지공시지가는 당해 공익사업에 따른 가격변동이 있는 것으로 조사됨 \par
주2) 기호 1, 3, 5, 7 표준지공시지가는 당해 공익사업과 직접 관계가 없는 것으로 조사됨 \par
주3) 기호 1, 2 표준지공시지가는 2013년, 기호 5, 8 표준지공시지가는 2021년 신규표준지임 \par
2. 인근 표준지공시지가의 연도별 공시가격 \par
1) 자연녹지지역 표준지공시지가 \par
(원/m²) \par
{\small\centering
\resizebox{\ifdim\width>\linewidth\linewidth\else\width\fi}{!}{%
\begin{tabular}{|c|c|c|c|c|c|c|c|}
\hline
\makecell{기호} & \makecell{소재지} & \makecell{2007.\\01.01.} & \makecell{2008.\\01.01.} & \makecell{2009.\\01.01.} & \makecell{2013.\\01.01.} & \makecell{2014.\\01.01.} & \makecell{2015.\\01.01.} \\
\hline
\makecell{1} & \makecell{A시 B구\\C동 60} & \makecell{-} & \makecell{-} & \makecell{-} & \makecell{70,000} & \makecell{75,000} & \makecell{85,000} \\
\hline
\makecell{2} & \makecell{A시 B구\\C동 80} & \makecell{-} & \makecell{-} & \makecell{-} & \makecell{45,000} & \makecell{50,000} & \makecell{55,000} \\
\hline
\makecell{3} & \makecell{A시 B구\\C동 120} & \makecell{42,000} & \makecell{44,000} & \makecell{44,000} & \makecell{50,000} & \makecell{55,000} & \makecell{60,000} \\
\hline
\makecell{4} & \makecell{A시 B구\\C동 140} & \makecell{100,000} & \makecell{102,000} & \makecell{102,000} & \makecell{110,000} & \makecell{115,000} & \makecell{120,000} \\
\hline
\end{tabular}}\par}

2) 주거지역 표준지공시지가 \par
(원/m²) \par
{\small\centering
\resizebox{\ifdim\width>\linewidth\linewidth\else\width\fi}{!}{%
\begin{tabular}{|c|c|c|c|c|c|c|c|}
\hline
\makecell{기호} & \makecell{소재지} & \makecell{2015.\\01.01.} & \makecell{2016.\\01.01.} & \makecell{2017.\\01.01.} & \makecell{2021.\\01.01.} & \makecell{2022.\\01.01.} & \makecell{2023.\\01.01.} \\
\hline
\makecell{5} & \makecell{A시 B구\\C동 160} & \makecell{-} & \makecell{-} & \makecell{-} & \makecell{450,000} & \makecell{460,000} & \makecell{460,000} \\
\hline
\makecell{6} & \makecell{A시 B구\\C동 180} & \makecell{167,000} & \makecell{174,000} & \makecell{175,000} & \makecell{240,000} & \makecell{250,000} & \makecell{230,000} \\
\hline
\makecell{7} & \makecell{A시 B구\\C동 220} & \makecell{143,000} & \makecell{146,000} & \makecell{147,000} & \makecell{185,000} & \makecell{189,000} & \makecell{173,000} \\
\hline
\makecell{8} & \makecell{A시 B구\\C동 240} & \makecell{-} & \makecell{-} & \makecell{-} & \makecell{310,000} & \makecell{320,000} & \makecell{300,000} \\
\hline
\end{tabular}}\par}

< 자료 5 > 지가변동률(A시 B구) \par
{\small\centering
\resizebox{\ifdim\width>\linewidth\linewidth\else\width\fi}{!}{%
\begin{tabular}{|c|c|c|}
\hline
\multirow{2}{*}{\makecell{기간}} & \multicolumn{2}{|c|}{\makecell{지가변동률(\%)}} \\
\cline{2-3}
 & \makecell{녹지지역} & \makecell{주거지역} \\
\hline
\makecell{2013.01.01.～2013.12.29.} & \makecell{6.00} & \makecell{5.00} \\
\hline
\makecell{2013.01.01.～2014.12.29.} & \makecell{10.00} & \makecell{8.00} \\
\hline
\makecell{2014.01.01.～2014.12.29.} & \makecell{4.00} & \makecell{3.00} \\
\hline
\makecell{2021.01.01.～2021.06.17.} & \makecell{5.00} & \makecell{4.00} \\
\hline
\makecell{2021.01.01.～2022.06.17.} & \makecell{8.00} & \makecell{6.00} \\
\hline
\makecell{2022.01.01.～2022.06.17.} & \makecell{3.00} & \makecell{2.00} \\
\hline
\end{tabular}}\par}

< 자료 6 > 지역요인, 개별요인 및 그 밖의 요인 \par
1. 본건 토지와 기호 1～8 표준지공시지가는 인근지역에 위치하므로 지역요인 대등함 \par
2. 개별요인 및 그 밖의 요인 비교 \par
{\small\centering
\resizebox{\ifdim\width>\linewidth\linewidth\else\width\fi}{!}{%
\begin{tabular}{|c|c|c|c|c|c|c|c|c|}
\hline
\makecell{구분} & \makecell{표준지\\1} & \makecell{표준지\\2} & \makecell{표준지\\3} & \makecell{표준지\\4} & \makecell{표준지\\5} & \makecell{표준지\\6} & \makecell{표준지\\7} & \makecell{표준지\\8} \\
\hline
\makecell{개별요인} & \makecell{1.150} & \makecell{1.100} & \makecell{1.050} & \makecell{1.050} & \makecell{1.150} & \makecell{1.080} & \makecell{1.050} & \makecell{1.050} \\
\hline
\makecell{그 밖의\\요인} & \makecell{2.00} & \makecell{2.00} & \makecell{2.00} & \makecell{2.00} & \makecell{3.00} & \makecell{3.00} & \makecell{3.00} & \makecell{3.00} \\
\hline
\end{tabular}}\par}

주1) 개별요인과 관련된 수치(= 평가대상 ÷ 표준지)는 평가대상과 각각의 표준지를 비교하여 산정된 결과임 \par
주2) 그 밖의 요인과 관련된 수치는 “표준지 기준 산정방식”에 의해 산정된 결과임 \par
주3) 개별요인 및 그 밖의 요인과 관련된 수치는 본건 일련번호 1, 2 토지에 동일하게 적용됨 \par
< 자료 7 > 기타 \par
1. 윤년은 고려치 않음 \par
2. 인근 표준지공시지가의 특성항목은 연도에 따른 변동은 없음 \par
3. 법원 확인 결과 일련번호 “1”토지는 2008.12.29.에, 일련번호 “2” 토지는 2016.06.17.에 협의취득에 따른 매매대금이 지급되었고, 같은 날 소유권이전등기를 마침 \par
4. 가격조사완료일: 2025.07.12. \par
5. 산정단가는 백원 단위에서 반올림함 \par
6. 인근 유사 토지의 지가변동률은 \%로 표시하되 소수점 첫째자리에서 반올림함 (예: 11.4\% 상승 -> 11\% 상승, 11.6\% 상승 -> 12\% 상승) \par
7. 상기자료 등은 법원, 감정평가사, 원고, 피고가 모두 확인하고 동의한 사항임 \par
\end{longquestion}
